\begin{papercard}{georgiou2026}
\cf{Stages} \stagetag{L} (two-model LLM extraction from policies with verbatim quotes, validated against manual labels), \stagetag{D} (policy, APK, permissions and traffic captured per app), \stagetag{A} (keyword search on Google Play for mental-health apps). arXiv preprint of May 2026 (UC Irvine, UC Riverside; De Cristofaro and Tsudik are established security and privacy researchers). Added from the gap search of 23 October 2026 for the mental-health scope asked for on 13 October and as a recent template for validated LLM extraction.

\cf{Sample} 25 popular Android apps for therapy and life coaching \pc{georgiou2026}{2}{this paper presents a multi-method measurement study of 25 popular Android apps offering mental health and life-coaching services}, found with five keywords \pc{georgiou2026}{2}{We identify candidate apps through Google Play searches using the following keywords: therapy, mental health, CBT, counseling, emotional wellness}; installs ranged from 10K+ to 10M+.

\cf{LLM extraction} Keyword matching first, then LLM extraction with a second model as cross-check \pc{georgiou2026}{2}{an ini- tial keyword-matching pass followed by LLM-assisted extrac- tion using Gemini Pro 3.1, validated against manual annota- tion and cross-checked with GPT-5.4}. The prompt asks for a binary label per category and the supporting verbatim quote \pc{georgiou2026}{5}{asks the model to return, for each of the five categories above, a binary label and supporting verbatim quotes from the policy text}. Validation: one author labelled the 24 unique policies \pc{georgiou2026}{5}{one author annotates the 24 unique privacy policies in our corpus (Headspace and Headspace Care share one) across the same five categories}, and every case where the two models disagreed was checked manually \pc{georgiou2026}{7}{We first validate extracted labels against manual annotation, including wherever the second model (GPT-5.4) disagrees}. Result: macro-F1 0.99 for Gemini and 0.97 for GPT-5.4 \pc{georgiou2026}{7}{Gemini Pro 3.1 achieves a macro-average F1 of 0.99, with precision 0.98 and recall 0.99, compared to GPT-5.4’s 0.97, 0.98, and 0.96, respectively}.

\cf{Findings} Every app embeds at least one tracker SDK its policy does not name, and 85\,\% fail to disclose at least half of them \pc{georgiou2026}{1}{every app in our corpus embeds at least one tracker SDK not named in its privacy policy, and 85\% of the apps we instrument fail to disclose at least half of the trackers detected in their APKs}. Twelve of 25 policies say data go to third-party AI providers; only five name the provider, seven use generic terms \pc{georgiou2026}{2}{The re- maining 7 use generic language, such as “AI Services” or “Large Language Models.”}. The authors conclude that disclosure falls short of informed consent \pc{georgiou2026}{1}{our results show that current disclosure practices fall short of the transparency needed for meaningful informed consent}.

\cf{Limits (card)} Preprint. Validation labels from one author, on the same 24 policies that were analysed, so it is a check, not a held-out test; five coarse binary categories are much easier than span-level harm coding. Android only, 25 apps, English policies.

\cf{Use for this thesis} (1) Mental-health apps (D3): a recent study that defines the domain by keywords and shows tracker and AI-provider disclosure gaps that the fitness and fertility harm trees can borrow. (2) D9 design pattern: require a verbatim quote for every positive label, run a second model, and send every disagreement to manual review; this is cheap and fits the screener role. (3) The validation has the thesis's limitation (one coder); report it the same way and do not read 0.99 as transferable to harder codes. (4) ``Generic language'' about recipients (``AI services'') is an Axis A/B coding example of vague disclosure.
\end{papercard}
